% GENERATED FILE. Edit canonical lessons, then run npm run generate:books.
% canonical-source-hash: fnv1a64:5022d0d7ef6a1947
% canonical-lessons: ZH-C340-chuantong, ZH-C340-jiangjiu, ZH-C340-guiju, ZH-C340-jieqi, ZH-C340-saomu, ZH-R340-review-336-340

\chapter{Old Ways}
\label{ch:zh-old-ways}
\hlchaptermodality{340}

\begin{chapteropening}
\textbf{By the end of this chapter:} \emph{I can talk about traditions, being particular about them, proper rules, the solar terms and tending graves.}

Complete the last lesson of chapter 340: I can talk about traditions, being particular about them, proper rules, the solar terms and tending graves.
\end{chapteropening}

\section[chuántǒng]{chuántǒng — tradition}

\label{lesson:ZH-C340-chuantong}



\begin{quote}
Before the new one: say the Mandarin for full of festive joy, then the Mandarin for fireworks lighting up the sky.
\end{quote}

\subsection*{You'll want to know: chuántǒng}
\emph{chuántǒng} — \textquotedblleft{}tradition\textquotedblright{}.

\begin{quote}
\emph{guònián chī jiǎozi shì yí ge lǎo chuántǒng} — Eating dumplings at the new year is an old tradition.
\end{quote}

Its characters come later; for now it is a word for the ear.

\begin{cousinweb}
\emph{chuántǒng} is built from \zh{传} 'to pass on' and \emph{tǒng}, both inherited from Middle Chinese, \emph{drjwen} and \emph{thowngH} (Guangyun, in Baxter's transcription), which glosses the second 'to gather together' and 'the end of a thread'.
\end{cousinweb}

\begin{grammarlens}[title={What you've built}]
One more word for telling what happened.
\end{grammarlens}

\subsection*{Your turn}
\begin{itemize}
\raggedright
  \item \emph{Say it:} \emph{chuántǒng}
  \item \emph{Say it:} \emph{chuántǒng}, once more
  \item \emph{Say it:} \emph{guònián chī jiǎozi shì yí ge lǎo chuántǒng} — Eating dumplings at the new year is an old tradition.
  \item \emph{Recall:} say the Mandarin for full of festive joy, then the Mandarin for fireworks lighting up the sky, then say \emph{chuántǒng} again
\end{itemize}

\begin{tcolorbox}[breakable,colback=teal!4,colframe=teal!35!black,title={Before you move on}]
What does chuántǒng mean? (\textquotedblleft{}Tradition\textquotedblright{}.) Say it once more.
\end{tcolorbox}
\section[jiǎngjiu]{jiǎngjiu — to be particular about}

\label{lesson:ZH-C340-jiangjiu}



\begin{quote}
Before the new one: say the Mandarin for fireworks lighting up the sky, then the Mandarin for an old tradition.
\end{quote}

\subsection*{You'll want to know: jiǎngjiu}
\emph{jiǎngjiu} — \textquotedblleft{}to be particular about\textquotedblright{}.

\begin{quote}
\emph{wǒ nǎinai hěn jiǎngjiu chī} — My grandma is very particular about food.
\end{quote}

Its characters come later; for now it is a word for the ear.

\begin{cousinweb}
\emph{jiǎngjiu} is built from \emph{jiǎng} and \emph{jiū}, both inherited from Middle Chinese, \emph{kaewngX} and \emph{kjuwH} (Guangyun, in Baxter's transcription), which glosses the first 'to discuss' and the second 'to go to the end' and 'deep'.
\end{cousinweb}

\begin{grammarlens}[title={What you've built}]
One more word for telling what happened.
\end{grammarlens}

\subsection*{Your turn}
\begin{itemize}
\raggedright
  \item \emph{Say it:} \emph{jiǎngjiu}
  \item \emph{Say it:} \emph{jiǎngjiu}, once more
  \item \emph{Say it:} \emph{wǒ nǎinai hěn jiǎngjiu chī} — My grandma is very particular about food.
  \item \emph{Recall:} say the Mandarin for fireworks lighting up the sky, then the Mandarin for an old tradition, then say \emph{jiǎngjiu} again
\end{itemize}

\begin{tcolorbox}[breakable,colback=teal!4,colframe=teal!35!black,title={Before you move on}]
What does jiǎngjiu mean? (\textquotedblleft{}To be particular about\textquotedblright{}.) Say it once more.
\end{tcolorbox}
\section[guīju]{guīju — rules of conduct; manners}

\label{lesson:ZH-C340-guiju}



\begin{quote}
Before the new one: say the Mandarin for an old tradition, then the Mandarin for to be particular about details.
\end{quote}

\subsection*{You'll want to know: guīju}
\emph{guīju} — \textquotedblleft{}rules of conduct; manners\textquotedblright{}.

\begin{quote}
\emph{chī niányèfàn yǒu shénme guīju} — What are the rules at the reunion dinner?
\end{quote}

Its characters come later; for now it is a word for the ear.

\begin{cousinweb}
\emph{guīju} is built from \emph{guī} 'a compass' and \emph{jǔ} 'a set square', both inherited from Middle Chinese, \emph{kjwie} and \emph{kjuX} (Guangyun, in Baxter's transcription), which glosses the first 'round', quoting an older book that uses this very pair, and the second 'a rule'.
\end{cousinweb}

\begin{grammarlens}[title={What you've built}]
One more word for telling what happened.
\end{grammarlens}

\subsection*{Your turn}
\begin{itemize}
\raggedright
  \item \emph{Say it:} \emph{guīju}
  \item \emph{Say it:} \emph{guīju}, once more
  \item \emph{Say it:} \emph{chī niányèfàn yǒu shénme guīju} — What are the rules at the reunion dinner?
  \item \emph{Recall:} say the Mandarin for an old tradition, then the Mandarin for to be particular about details, then say \emph{guīju} again
\end{itemize}

\begin{tcolorbox}[breakable,colback=teal!4,colframe=teal!35!black,title={Before you move on}]
What does guīju mean? (\textquotedblleft{}Rules of conduct; manners\textquotedblright{}.) Say it once more.
\end{tcolorbox}
\section[\texorpdfstring{\zh{节气} (jiéqì)}{节气 (jiéqì)}]{\zh{节气} (jiéqì) — solar term}

\label{lesson:ZH-C340-jieqi}



\begin{quote}
Before the new one: say the Mandarin for to be particular about details, then the Mandarin for the rules of conduct.
\end{quote}

\subsection*{You'll want to know: \zh{节气}}
\textbf{\zh{节气}} — \emph{jiéqì} — \textquotedblleft{}solar term\textquotedblright{}.

\begin{quote}
\textbf{\zh{每年有二十四个节气}} \emph{měinián yǒu èrshísì ge jiéqì} — Every year has twenty-four solar terms.
\end{quote}

\begin{cousinweb}
\textbf{\zh{节气}} (\emph{jiéqì}) is built from \zh{节} 'a joint' and \zh{气} 'breath', both inherited from Middle Chinese, \emph{tset} and \emph{khj+jH} (Guangyun, in Baxter's transcription), which glosses the first 'a joint of bamboo' and the second 'breath'. The calendar sense is not given there.
\end{cousinweb}

\begin{grammarlens}[title={What you've built}]
One more word for telling what happened.
\end{grammarlens}

\subsection*{Your turn}
\begin{itemize}
\raggedright
  \item \emph{Say it:} \emph{jiéqì}
  \item \emph{Say it:} \emph{jiéqì}, once more
  \item \emph{Say it:} \emph{měinián yǒu èrshísì ge jiéqì} — Every year has twenty-four solar terms.
  \item \emph{Recall:} say the Mandarin for to be particular about details, then the Mandarin for the rules of conduct, then say \emph{jiéqì} again
\end{itemize}

\begin{tcolorbox}[breakable,colback=teal!4,colframe=teal!35!black,title={Before you move on}]
What does \zh{节气} mean? (\textquotedblleft{}Solar term\textquotedblright{}.) Say it once more.
\end{tcolorbox}
\section[sǎomù]{sǎomù — to tend the family graves}

\label{lesson:ZH-C340-saomu}



\begin{quote}
Before the new one: say the Mandarin for the rules of conduct, then the Mandarin for a solar term of the year.
\end{quote}

\subsection*{You'll want to know: sǎomù}
\emph{sǎomù} — \textquotedblleft{}to tend the family graves\textquotedblright{}.

\begin{quote}
\emph{měinián chūntiān wǒmen dōu qù sǎomù} — Every spring we go to tend the family graves.
\end{quote}

Its characters come later; for now it is a word for the ear.

\begin{cousinweb}
\emph{sǎomù} is built from \zh{扫} 'to sweep' and \emph{mù} 'a grave', both inherited from Middle Chinese, \emph{sawX} and \emph{muH} (Guangyun, in Baxter's transcription), which glosses the first 'to sweep clean' and the second 'a grave mound'.
\end{cousinweb}

\begin{grammarlens}[title={What you've built}]
One more word for telling what happened.
\end{grammarlens}

\subsection*{Your turn}
\begin{itemize}
\raggedright
  \item \emph{Say it:} \emph{sǎomù}
  \item \emph{Say it:} \emph{sǎomù}, once more
  \item \emph{Say it:} \emph{měinián chūntiān wǒmen dōu qù sǎomù} — Every spring we go to tend the family graves.
  \item \emph{Recall:} say the Mandarin for the rules of conduct, then the Mandarin for a solar term of the year, then say \emph{sǎomù} again
\end{itemize}

\begin{tcolorbox}[breakable,colback=teal!4,colframe=teal!35!black,title={Before you move on}]
What does sǎomù mean? (\textquotedblleft{}To tend the family graves\textquotedblright{}.) Say it once more.
\end{tcolorbox}
\section[(dialogue)]{Looking back: from Seeing Friends Off to Old Ways}

\label{lesson:ZH-R340-review-336-340}



\begin{quote}
Say the Mandarin for to take leave of someone and to give a send-off.
\end{quote}

\subsection*{Your turn}
\begin{itemize}
\raggedright
  \item \emph{Say it:} to take leave of someone, to give a send-off, to see a guest out, to give a present, to shake hands with someone
  \item \emph{Say it:} to congratulate someone, to wish someone well, a token of regard, good luck, to make a wish
  \item \emph{Say it:} to celebrate a festival, to pay a visit to someone, to be reunited as a family, a day off work, a local custom
  \item \emph{Say it:} a long holiday, the National Day holiday, to celebrate a happy event, full of festive joy, fireworks lighting up the sky
  \item \emph{Say it:} an old tradition, to be particular about details, the rules of conduct, a solar term of the year, to tend the family graves
  \item \emph{Read it:} \zh{送} \zh{祝} \zh{庆} — each one aloud, then a word that holds it
\end{itemize}

\begin{tcolorbox}[breakable,colback=teal!4,colframe=teal!35!black,title={Before you move on}]
To take leave of someone? (\emph{gàobié}.) To give a send-off? (\textbf{\zh{欢送}}, \emph{huānsòng}.) To see a guest out? (\textbf{\zh{送客}}, \emph{sòngkè}.) To give a present? (\emph{sònglǐ}.) To shake hands with someone? (\emph{wòshǒu}.) To congratulate someone? (\emph{gōngxǐ}.) To wish someone well? (\textbf{\zh{祝}}, \emph{zhù}.) A token of regard? (\textbf{\zh{心意}}, \emph{xīnyì}.) Good luck? (\emph{hǎoyùn}.) To make a wish? (\emph{xǔyuàn}.) To celebrate a festival? (\textbf{\zh{过节}}, \emph{guòjié}.) To pay a visit to someone? (\emph{bàifǎng}.) To be reunited as a family? (\emph{tuányuán}.) A day off work? (\textbf{\zh{假日}}, \emph{jiàrì}.) A local custom? (\emph{fēngsú}.) A long holiday? (\textbf{\zh{长假}}, \emph{chángjià}.) The National Day holiday? (\textbf{\zh{国庆}}, \emph{Guóqìng}.) To celebrate a happy event? (\textbf{\zh{庆祝}}, \emph{qìngzhù}.) Full of festive joy? (\textbf{\zh{喜庆}}, \emph{xǐqìng}.) Fireworks lighting up the sky? (\emph{yānhuā}.) An old tradition? (\emph{chuántǒng}.) To be particular about details? (\emph{jiǎngjiu}.) The rules of conduct? (\emph{guīju}.) A solar term of the year? (\textbf{\zh{节气}}, \emph{jiéqì}.) To tend the family graves? (\emph{sǎomù}.)
\end{tcolorbox}
